\section{Conclusiones}

El desarrollo de esta práctica se llevó a cabo gracias al tutorial de Raspberry Pi de la facultad \cite{FI_UNAM_Raspberry_2026} y a las clases del profesor Moisés \cite{Melendez_Reyes_2026}.

\subsection{Hernandez Acosta Mauricio Gabriel}
Esta práctica me ayudó mucho a entender cómo un microcontrolador realmente interactúa con lo que pasa en el mundo exterior. A veces uno se acostumbra a trabajar únicamente con señales digitales de ceros y unos, por lo que ver cómo el ADC traduce una variación continua de voltaje a valores numéricos en Thonny me pareció bastante intuitivo y revelador. Me gustó mucho el ejercicio del PWM con desvanecimiento, ya que pude comprobar en físico cómo al cambiar el ciclo de trabajo en fracciones de segundo la vista percibe un cambio gradual de brillo en el LED en lugar de parpadeos bruscos. En general, la práctica me dejó un aprendizaje muy claro sobre cómo leer sensores analógicos y cómo usar MicroPython para manipular periféricos de salida sin tantas complicaciones de bajo nivel.

\subsection{Perez Olvera Alexis Abraham}
Lo que más me gustó de esta práctica fue la parte de comparar los dos sensores de temperatura. Al principio pensé que el sensor interno del RP2040 sería suficiente para cualquier medición, pero al poner a prueba el TMP36 nos dimos cuenta de que el chip se calienta un poco por su propia operación y eso altera las lecturas, algo que en la teoría se menciona pero que no te queda del todo claro hasta que lo ves corriendo en la consola. También me pareció muy satisfactorio el último ejercicio, donde conectamos el potenciómetro para gobernar las dos salidas PWM complementarias; girar la perilla y ver cómo la intensidad de un LED subía mientras el otro se apagaba de forma suave y sin retrasos demostró lo útil que es combinar una entrada analógica con salidas moduladas para tareas de control en tiempo real.

\subsection{Sierra Garcia Mariana}
Personalmente encontré muy entretenida y formativa esta práctica, sobre todo al trabajar con la fotoresistencia y armar el divisor de voltaje. Fue muy gratificante tapar el sensor con la mano y ver cómo el programa detectaba la caída de luz en tiempo real para activar la salida en el GPIO7, funcionando justo como el encendido automático de las luces públicas. Además, programar las señales PWM inversas en los pines 0 y 1 me ayudó a entender mejor la lógica detrás de los actuadores complementarios, como los que se usan para cambiar el sentido de giro de un motor o en controles de potencia balanceados. Siento que esta sesión nos dio herramientas muy prácticas que sin duda aplicaremos en proyectos más avanzados del laboratorio.
